\assignment{5}{Mon 03 Oct 2022 20:09}{}

\begin{enumerate}
	\item An electron is trapped in a one-dimensional region of width $0.062 \mathrm{~nm}$. Find the three smallest possible values allowed for the energy of the electron. 

	\item An electron is trapped in a one-dimensional well of width $0.132 \mathrm{~nm}$. The electron is in the $n=10$ state. (a) What is the energy of the electron? (b) What is the uncertainty in its momentum? (c) What is the uncertainty in its position? (d) How do these results change as $n \rightarrow \infty$ ? Is this consistent with classical behavior?

\item A proton or a neutron can sometimes "violate" conservation of energy by emitting and then reabsorbing a pi meson, which has a mass of $135 \mathrm{MeV} / \mathrm{c}^2$. This is possible as long as the pi meson is reabsorbed within a short enough time $\Delta t$ consistent with the uncertainty principle. A) Consider $p \rightarrow p+\pi$. But what amount $\Delta E$ is the energy conservation violated? (Ignore any kinetic energies.) B) For how long a time $\Delta t$ can the pi meson exist? C) Assuming that the pi meson travels at very nearly the speed of light, how far from the proton can it go?
\end{enumerate}
